\documentclass[14pt]{article}
\input{preamble}

\begin{document}

\begin{titlepage}
\begin{center}

Министерство цифрового развития, связи и массовых\\
коммуникаций Российской Федерации\\[4pt]
Федеральное государственное бюджетное общеобразовательное\\
учреждение высшего образования\\[4pt]
<<Сибирский государственный университет телекоммуникаций\\
и информатики>>\\[4pt]
(СибГУТИ)
\vspace{1cm}

Факультет комплексной безопасности ТЭК\\
Кафедра информатики\\

\vspace{2cm}

\textbf{\Large ОТЧЁТ}\\
\textbf{\large по дисциплине <<Интернет технологии и сетовое ПО>>}\\

\vspace{0.5cm}

\textbf{\large Неделя 17 --- Финальный проект}\\
\textbf{\large <<Микросервисная архитектура интернет-магазина>>}

\vspace{3cm}

\end{center}

\begin{flushright}
\begin{minipage}{0.45\textwidth}
Выполнил:\\
студент группы ИКС-433\\
Рудометов Д.Р.\\
Проверил:\\
Лошкарев А.В.
\end{minipage}
\end{flushright}

\vfill

\begin{center}
Новосибирск, 2026
\end{center}

\end{titlepage}


\tableofcontents
\newpage

\section{Введение}

Целью данной работы является проектирование и реализация микросервисной архитектуры для интернет-магазина. Проект объединяет знания, полученные на протяжении курса: REST API, gRPC, Docker, очереди сообщений (RabbitMQ), базы данных PostgreSQL.

Система состоит из нескольких независимых сервисов, каждый из которых отвечает за свою предметную область: каталог товаров, обработка заказов, уведомления и единая точка входа (API Gateway).

\section{Постановка задачи}

Необходимо:
\begin{enumerate}
    \item Спроектировать систему из нескольких микросервисов.
    \item Описать архитектуру в \texttt{ARCHITECTURE.md}.
    \item Реализовать код сервисов с использованием REST и gRPC.
    \item Упаковать всё в Docker и обеспечить запуск через \texttt{docker-compose up}.
    \item Настроить CI/CD конвейер (GitHub Actions) для автоматической проверки, сборки и деплоя.
    \item Подготовить Helm-чарты для развертывания в Kubernetes.
    \item Обеспечить слабую связность: падение одного сервиса не должно блокировать другие.
\end{enumerate}

\section{Архитектура системы}

\subsection{Обзор компонентов}

Проект состоит из следующих сервисов:

\begin{itemize}
    \item \textbf{API Gateway} --- единая точка входа, проксирует запросы к внутренним сервисам по HTTP.
    \item \textbf{Catalog Service} --- управление каталогом товаров (REST API + gRPC сервер).
    \item \textbf{Order Service} --- создание и просмотр заказов, проверка товара через gRPC, отправка событий в RabbitMQ.
    \item \textbf{Notification Service} --- потребитель очереди, обрабатывает события о создании заказов.
    \item \textbf{PostgreSQL} --- реляционная база данных для хранения товаров и заказов.
    \item \textbf{RabbitMQ} --- брокер сообщений для асинхронного взаимодействия между сервисами.
\end{itemize}

\subsection{Схема взаимодействия}

Внешние клиенты отправляют HTTP-запросы к API Gateway (порт 8000), который проксирует их к соответствующим сервисам:

\begin{itemize}
    \item \texttt{/api/products} $\rightarrow$ Catalog Service (HTTP)
    \item \texttt{/api/orders} $\rightarrow$ Order Service (HTTP)
\end{itemize}

При создании заказа Order Service:
\begin{enumerate}
    \item Проверяет наличие товара через \textbf{gRPC} вызов к Catalog Service (порт 50051).
    \item Рассчитывает стоимость заказа.
    \item Сохраняет заказ в PostgreSQL.
    \item Публикует событие \texttt{OrderCreated} в очередь \texttt{notifications} через \textbf{RabbitMQ}.
\end{enumerate}

Notification Service прослушивает очередь и логирует уведомления.

\subsection{Используемые протоколы}

\begin{center}
\begin{tabular}{|l|l|l|}
\hline
\textbf{Связь} & \textbf{Протокол} & \textbf{Обоснование} \\
\hline
Клиент $\rightarrow$ API Gateway & HTTP/REST & Универсальность, удобство \\
\hline
API Gateway $\rightarrow$ Сервисы & HTTP/REST & Проксирование запросов \\
\hline
Order $\rightarrow$ Catalog & gRPC & Низкая задержка, типизация \\
\hline
Order $\rightarrow$ Notification & RabbitMQ & Асинхронность, надёжность \\
\hline
\end{tabular}
\end{center}

\section{Реализация}

\subsection{Стек технологий}

\begin{itemize}
    \item \textbf{Язык}: Python 3.10
    \item \textbf{Фреймворк}: FastAPI (асинхронный)
    \item \textbf{ORM}: SQLAlchemy (asyncpg)
    \item \textbf{gRPC}: grpcio, grpc-tools
    \item \textbf{Брокер}: RabbitMQ 3 (aio-pika)
    \item \textbf{БД}: PostgreSQL 15
    \item \textbf{Контейнеризация}: Docker, Docker Compose
\end{itemize}

\subsection{Catalog Service}

Сервис каталога предоставляет REST API для CRUD-операций с товарами и gRPC-интерфейс для проверки наличия товара.

\textbf{Модель данных (models.py):}
\begin{lstlisting}[language=Python]
class Product(Base):
    __tablename__ = "products"
    id = Column(Integer, primary_key=True, index=True)
    name = Column(String, index=True)
    description = Column(String, nullable=True)
    price = Column(Float, nullable=False)
    stock = Column(Integer, default=0)
\end{lstlisting}

\textbf{Proto-файл (catalog.proto):}
\begin{lstlisting}[language=C]
syntax = "proto3";
package catalog;

service CatalogService {
  rpc CheckProduct (CheckProductRequest)
      returns (CheckProductResponse);
}

message CheckProductRequest {
  int32 product_id = 1;
}

message CheckProductResponse {
  bool exists = 1;
  double price = 2;
  string name = 3;
}
\end{lstlisting}

\textbf{REST-эндпоинты:}
\begin{itemize}
    \item \texttt{POST /api/products} --- создание товара
    \item \texttt{GET /api/products} --- список товаров
\end{itemize}

\subsection{Order Service}

Сервис заказов использует gRPC для проверки товара в каталоге и RabbitMQ для отправки уведомлений.

\textbf{Модель данных (models.py):}
\begin{lstlisting}[language=Python]
class Order(Base):
    __tablename__ = "orders"
    id = Column(Integer, primary_key=True, index=True)
    product_id = Column(Integer, index=True)
    quantity = Column(Integer, nullable=False)
    total_price = Column(Float, nullable=False)
    status = Column(String, default="CREATED")
    customer_email = Column(String)
\end{lstlisting}

\textbf{Ключевая логика создания заказа:}
\begin{lstlisting}[language=Python]
@app.post("/api/orders", response_model=OrderOut)
async def create_order(order: OrderCreate,
                       session: AsyncSession = Depends(get_session)):
    catalog_response = await check_product_in_catalog(
        order.product_id)
    if not catalog_response.exists:
        raise HTTPException(status_code=404,
            detail="Product not found in catalog")

    total_price = catalog_response.price * order.quantity

    new_order = Order(
        product_id=order.product_id,
        quantity=order.quantity,
        total_price=total_price,
        customer_email=order.customer_email,
        status="CREATED"
    )
    session.add(new_order)
    await session.commit()

    if mq_channel:
        message_body = json.dumps({
            "event": "OrderCreated",
            "order_id": new_order.id,
            "product_name": catalog_response.name,
            "total_price": new_order.total_price,
            "customer_email": new_order.customer_email
        }).encode()
        await mq_channel.default_exchange.publish(
            aio_pika.Message(body=message_body),
            routing_key="notifications",
        )
    return new_order
\end{lstlisting}

\subsection{API Gateway}

API Gateway реализован как обратный прокси на базе FastAPI и httpx. Он маршрутизирует запросы:
\begin{itemize}
    \item \texttt{/api/products*} $\rightarrow$ \texttt{catalog-svc:8000}
    \item \texttt{/api/orders*} $\rightarrow$ \texttt{order-svc:8000}
\end{itemize}

\begin{lstlisting}[language=Python]
async def proxy_request(request, base_url, path):
    url = f"{base_url}/{path}"
    async with httpx.AsyncClient() as client:
        req_body = await request.body()
        response = await client.request(
            method=request.method,
            url=url,
            headers=dict(request.headers),
            content=req_body,
            params=request.query_params
        )
        return JSONResponse(
            content=response.json(),
            status_code=response.status_code)
\end{lstlisting}

Также предоставляет эндпоинт \texttt{/health} для проверки состояния.

\subsection{Notification Service}

Сервис уведомлений --- автономный консьюмер RabbitMQ, который обрабатывает события из очереди \texttt{notifications}:

\begin{lstlisting}[language=Python]
async def process_message(message):
    async with message.process():
        body = json.loads(message.body.decode())
        event = body.get("event")
        if event == "OrderCreated":
            logger.info(
                f"SENDING EMAIL to {body['customer_email']}: "
                f"Order #{body['order_id']} created!")
\end{lstlisting}

Сервис устойчив к недоступности RabbitMQ: в цикле ожидает подключения с паузой 2 секунды.

\section{Инфраструктура}

\subsection{Docker}

Каждый сервис упакован в Docker-образ на базе \texttt{python:3.10-slim}. Dockerfile для Catalog Service включает сборку gRPC-кода из proto-файла:

\begin{lstlisting}[language=bash]
FROM python:3.10-slim
WORKDIR /app
RUN apt-get update && apt-get install -y gcc g++
COPY requirements.txt .
RUN pip install --no-cache-dir -r requirements.txt
COPY . .
RUN python -m grpc_tools.protoc -I. \
    --python_out=. --grpc_python_out=. catalog.proto
CMD ["python", "main.py"]
\end{lstlisting}

\subsection{Docker Compose}

Файл \texttt{docker-compose.yml} описывает все 6 сервисов:

\begin{lstlisting}[language=yaml]
version: '3.8'
services:
  postgres:
    image: postgres:15-alpine
    environment:
      POSTGRES_USER: user
      POSTGRES_PASSWORD: password
      POSTGRES_DB: shop_db
    healthcheck:
      test: ["CMD-SHELL",
             "pg_isready -U user -d shop_db"]

  rabbitmq:
    image: rabbitmq:3-management-alpine
    healthcheck:
      test: rabbitmq-diagnostics -q ping

  api-gateway:
    build: ./api-gateway
    ports: ["8000:8000"]
    depends_on: [catalog-svc, order-svc]

  catalog-svc:
    build: ./catalog-svc
    depends_on:
      postgres: {condition: service_healthy}

  order-svc:
    build: ./order-svc
    depends_on:
      postgres: {condition: service_healthy}
      rabbitmq: {condition: service_healthy}

  notification-svc:
    build: ./notification-svc
    depends_on:
      rabbitmq: {condition: service_healthy}
\end{lstlisting}

Используются healthcheck'и для PostgreSQL и RabbitMQ, чтобы зависимые сервисы запускались только после готовности инфраструктуры.

\subsection{Запуск проекта}

\begin{lstlisting}[language=bash]
docker-compose up --build
\end{lstlisting}

После запуска API доступен на \texttt{http://localhost:8000}.

\subsection{Helm и Kubernetes}

Для развертывания приложения в кластере Kubernetes подготовлен Helm-чарт (\texttt{helm/shop}). Он включает в себя манифесты Deployment и Service для всех микросервисов, а также зависимости для базы данных и брокера сообщений. Чарт поддерживает переопределение параметров для различных окружений (файлы \texttt{values-staging.yaml} и \texttt{values-production.yaml}).

\subsection{CI/CD Пайплайн (GitHub Actions)}

Настроен полноценный конвейер непрерывной интеграции и доставки (CI/CD) с использованием GitHub Actions. Пайплайн состоит из 6 стадий:
\begin{enumerate}
    \item \textbf{Lint}: Проверка кода Python (\texttt{flake8}), Docker-файлов (\texttt{hadolint}) и Helm-чартов (\texttt{helm lint}).
    \item \textbf{Build \& Test}: Сборка Docker-образов и запуск интеграционных тестов.
    \item \textbf{Security Scan}: Поиск уязвимостей в собранных образах с помощью \texttt{Trivy}.
    \item \textbf{Publish}: Публикация протестированных образов в GitHub Container Registry (ghcr.io).
    \item \textbf{Deploy Staging}: Генерация манифестов Helm и деплой в staging окружение.
    \item \textbf{Deploy Production}: Деплой в production окружение (с требованием ручного подтверждения).
\end{enumerate}

\section{Связность и отказоустойчивость}

\begin{itemize}
    \item \textbf{Слабая связность}: сервисы взаимодействуют через сетевые протоколы (HTTP, gRPC, AMQP), а не через общий код. Падение Notification Service не влияет на создание заказов.
    \item \textbf{Graceful degradation}: Order Service корректно обрабатывает недоступность RabbitMQ --- заказ создаётся, но уведомление не отправляется.
    \item \textbf{Retry-логика}: Notification Service повторяет попытки подключения к RabbitMQ с интервалом 2 секунды.
    \item \textbf{Healthcheck}: PostgreSQL и RabbitMQ проверяются на готовность перед запуском зависимых сервисов.
\end{itemize}

\section{Ответы на контрольные вопросы}

\textbf{1. Какие критерии использовались при разделении на микросервисы?}

Разделение выполнено по принципу \emph{bounded context} (DDD): каждый сервис отвечает за свою предметную область --- каталог, заказы, уведомления. Критерии: независимость развёртывания, единая ответственность, минимизация связей.

\textbf{2. Почему выбран именно этот стек технологий?}

Python и FastAPI обеспечивают высокую скорость разработки и встроенную поддержку async/await. PostgreSQL~--- надёжная СУБД с поддержкой транзакций. gRPC~--- эффективный протокол для межсервисного взаимодействия. RabbitMQ~--- зрелый брокер сообщений.

\textbf{3. Как сервисы справляются с ошибками?}

Order Service перехватывает \texttt{grpc.RpcError} и возвращает HTTP 503 при недоступности Catalog. Подключение к RabbitMQ происходит через \texttt{connect\_robust} с автопереподключением. Notification Service использует цикл с \texttt{await asyncio.sleep(2)} для ожидания RabbitMQ.

\textbf{4. Как организован деплой?}

Локально~--- через \texttt{docker-compose up}. Продакшен деплой автоматизирован через 6-стадийный CI/CD пайплайн GitHub Actions, который собирает образы, проверяет их на уязвимости, публикует в ghcr.io и разворачивает в Kubernetes с помощью Helm-чартов (окружения staging и production).

\textbf{5. Что было самым сложным?}

Обеспечение корректного порядка запуска сервисов (healthcheck и depends\_on) и интеграция gRPC-сервера в lifespan FastAPI-приложения.

\textbf{6. Что можно улучшить?}

Добавить распределённый трейсинг (OpenTelemetry), централизованный сбор логов (ELK), circuit breaker для gRPC-вызовов, и полноценную отправку email.

\section{Заключение}

В ходе выполнения финального проекта была спроектирована и реализована микросервисная система интернет-магазина, включающая:
\begin{itemize}
    \item 4 микросервиса с чётким разделением ответственности;
    \item межсервисное взаимодействие через REST, gRPC и RabbitMQ;
    \item полную контейнеризацию с Docker Compose;
    \item автоматизированный CI/CD пайплайн (GitHub Actions) и Helm-чарты для Kubernetes;
    \item обработку ошибок и graceful degradation.
\end{itemize}

Проект запускается одной командой \texttt{docker-compose up} и демонстрирует применение ключевых паттернов распределённых систем.

\end{document}